\emph{Setting.} The module, the generators, the structure $h$, the null vector $X(x)$ and the incidence line as the vacua of $X$ are those of Theorem \claimref{synthesis.incidence-vacua-thermofield-real-capacity} and Proposition \claimref{one-relation.spacetime-manifold-vacuum-lines}, extended complex-linearly to complex points $m$; $Q(m) = \sum_{a} m_{a}^{2}Q(e_{a})$ is the complex-bilinear extension, and $X(m)$ is null for every $m$: $\nO$ and $\ninf$ are null and orthogonal to $m$, and $\nO\ninf + \ninf\nO = 2$ (\claimref{std-algebra.vector-squares}), so $X(m)^{2} = Q(m) - \tfrac12 Q(m)(\nO\ninf + \ninf\nO) = 0$. The check realises the module as three fermionic modes in the Jordan-Wigner basis, with the identity Gram matrix as $h$, the Hermitian Majorana operators as the generators of square $+1$ and $i$ times them as the generators of square $-1$.

\emph{(i)} If $h_{1}, h_{2}$ both make every generator unitary then $h_{1}^{-1}h_{2}$ commutes with the generators, which generate the full matrix algebra of the irreducible module, so it is a scalar; the check confirms that the Hermitian forms for which the six generators are unitary form one real line. $J^{2} = i^{2}e_{1}e_{+}e_{1}e_{+} = -(-e_{1}^{2}e_{+}^{2}) = 1$. $J^{\dagger} = -i\,e_{+}^{\dagger}e_{1}^{\dagger} = -i\,e_{+}e_{1} = i\,e_{1}e_{+} = J$, because $e_{1}, e_{+}$ square to $+1$ and are unitary, hence Hermitian. $J$ is a bivector times $i$, so even, and preserves each chirality half. $\Sigma$ is Hermitian because $J$ is $h$-Hermitian. For a bivector $X = g_{a}g_{b}$, $X^{\dagger} = \epsilon_{a}\epsilon_{b}\,g_{b}g_{a}$ with $\epsilon = +1$ for the Hermitian generators $e_{1}, e_{+}$ and $-1$ for the others, and $J$ commutes with $g_{a}g_{b}$ when $a, b$ are both in $\{e_{1}, e_{+}\}$ or both outside it and anticommutes otherwise; in each case $X^{\dagger}J + JX = 0$, which is $\Spin(2,4)$-invariance of $\Sigma$ (the check verifies all fifteen). On a half $J$ has eigenvalues $+1, +1, -1, -1$ (check), and $h$ is positive, so $\Sigma$ has signature $(2,2)$ there. Uniqueness: an invariant Hermitian form on the irreducible half $\CS$ corresponds to an intertwiner from $\CS$ to the conjugate dual; by Schur these form a space of complex dimension at most one, whose Hermitian elements are a real line. The check solves the linear system for Hermitian $M$ on the half with $R^{\dagger}M + MR = 0$ for the fifteen bivectors and finds a real solution space of dimension one. Proposition \claimref{correspondence.reality-conjugation-twistor-form-factor}(2) and (4) give the twistor form as a $\Spin(2,4)$-invariant Hermitian form of signature $(2,2)$ on $\CS$, so it is a nonzero real multiple of $\Sigma$.

\emph{(ii)} $X(ie_{1}) = \nO + ie_{1} - \tfrac12 Q(ie_{1})\ninf$ with $Q(ie_{1}) = -1$, so $X(ie_{1}) = \tfrac12(e_{+} + e_{-}) + ie_{1} + \tfrac12(e_{+} - e_{-}) = e_{+} + ie_{1}$, and $e_{+}(1 - J) = e_{+} - ie_{+}e_{1}e_{+} = e_{+} + ie_{1}e_{+}^{2} = e_{+} + ie_{1}$. Likewise $X(-ie_{1}) = e_{+} - ie_{1} = e_{+}(1 + J)$. Since $e_{+}$ is invertible, the vacua of $X(\pm ie_{1})$ are the $\pm 1$ eigenvectors of $J$, so on a half $L(ie_{1})$ and $L(-ie_{1})$ are the two eigenplanes of $J$; they are $h$-orthogonal and $\Sigma$-orthogonal (eigenspaces of an $h$-Hermitian involution), and $\Sigma = h(\cdot, J\cdot)$ equals $h$ on the first and $-h$ on the second. The check computes both lines as kernels and confirms each statement.

\emph{(iii)} For a generator $g$, $g^{\sharp} = J g^{\dagger}J$ (the $\Sigma$-adjoint, as $J^{-1} = J$ and $J^{\dagger} = J$). The check finds $g^{\sharp} = -g$ for all six; this is the $\Sigma$-anti-Hermiticity of real vectors, of which (i) is the bivector consequence. The adjoint is antilinear in coefficients, so $X(m)^{\sharp} = -X(\bar m)$, the coefficients of $\nO$ and $\ninf$ being real and $\overline{Q(m)} = Q(\bar m)$. A nonzero null $X$ acting on the eight-dimensional module has $X^{2} = 0$ and rank four (check, for every point used), so its kernel equals its image, and $X$ is odd. Hence a vacuum $\psi \in L(m)$ is $\psi = X(m)\chi$ with $\chi$ in the other half, and for $\varphi \in L(\bar m)$, $\Sigma(\psi, \varphi) = \Sigma(\chi, X(m)^{\sharp}\varphi) = -\Sigma(\chi, X(\bar m)\varphi) = 0$. By (i) and Theorem \claimref{synthesis.twistor-form-states-only-off-the-real-slice}(ii), read for $\Sigma$ under the basis map sending that theorem's time direction $e_{0}$ to $e_{1}$ here and its spatial $e_{1}, e_{2}, e_{3}$ to $e_{2}, e_{3}, e_{4}$, $L(m)$ is $\Sigma$-definite exactly when $Q(y) > 0$, positive for one time orientation of $y$; by (ii), $L(ie_{1})$ is positive, which fixes that orientation as $y^{1} > 0$ (the check confirms positivity at four future points, negativity at a past point, a split line at a spacelike imaginary part and a null line at a real point, and at each point computes the basis-free ratio $\det(\Sigma|_{L(m)})/\det(h|_{L(m)}) = 4Q(y)/\bigl(|1 - Q(m)|^{2} + 4|m^{1}|^{2} - 4Q(y)\bigr)$, with $m^{1}$ the $e_{1}$-component of $m$; this formula is computed at those points, not proved in general, and is not part of the statement). Then $L(\bar m)$ is negative ($\bar m$ has imaginary part $-y$), $L(m) \cap L(\bar m) = 0$ (a vector in both would be $\Sigma$-null in a definite plane), so $\CS = L(m) \oplus L(\bar m)$ $\Sigma$-orthogonally, and $J_{m} = \pm 1$ on $L(m)$, $L(\bar m)$ is an involution with $h_{m} = \Sigma(\cdot, J_{m}\cdot)$ positive. Conversely, let $h'$ be positive with $J' := h'^{-1}\Sigma$ an involution. $J'$ is $h'$-Hermitian (both forms are Hermitian), so its eigenplanes $L_{\pm}$ are $h'$-orthogonal with $\Sigma = \pm h'$ on them: $L_{+}$ is a $\Sigma$-positive plane (dimension two by the signature) and $L_{-}$ its $\Sigma$-orthogonal complement. $\Som$, the line of the point at infinity (in the model, the vacua of $\ninf$ in the half; the check confirms it is a $\Sigma$-null plane skew to every affine line it computes), is $\Sigma$-null (Proposition \claimref{correspondence.reality-conjugation-twistor-form-factor}(3)), so $L_{+} \cap \Som = 0$ and $L_{+}$ projects isomorphically onto $\Spi$: it is a graph line $\{A\pi + \pi\}$, the incidence line of exactly one point $m$ of the complexified affine chart (Theorem \claimref{correspondence.incidence-relation-chart}). Being positive, $m$ has future timelike imaginary part, and then $L_{-} = L_{+}^{\perp} = L(\bar m)$ by the orthogonality just proved and dimension. So $h' = h_{m}$, and by (ii) $h = h_{ie_{1}}$.

\emph{(iv)} If $s$ is definite, normalised trace against the Gram matrix gives a state. Suppose $s$ is indefinite. Take an $s$-orthogonal basis $f_{1}, \dots, f_{n}$ with $s(f_{i}, f_{i}) = \varepsilon_{i} \ne 0$, both signs occurring, and the matrix units $E_{ij}f_{k} = \delta_{jk}f_{i}$. Their adjoint is $E_{ij}^{\sharp} = (\varepsilon_{i}/\varepsilon_{j})E_{ji}$, so $E_{ij}^{\sharp}E_{ij} = (\varepsilon_{i}/\varepsilon_{j})E_{jj}$. For a positive $\omega$, choosing $i$ with $\varepsilon_{i}/\varepsilon_{j} > 0$ gives $\omega(E_{jj}) \ge 0$ and choosing $i$ with $\varepsilon_{i}/\varepsilon_{j} < 0$ gives $\omega(E_{jj}) \le 0$; so $\omega(E_{jj}) = 0$ for every $j$ and $\omega(1) = \sum_{j}\omega(E_{jj}) = 0$. The check builds a $\Sigma$-orthogonal basis of a half from the two lines of (ii), finds both signs, and verifies $E_{ij}^{\sharp}E_{ij} = (\varepsilon_{i}/\varepsilon_{j})E_{jj}$ for all sixteen pairs and that both signs multiply each $E_{jj}$. The complexified Clifford algebra is the full matrix algebra of the module, where $\Sigma$ has signature $(4,4)$, so the same argument applies. A state is therefore positive for some positive definite form, for which every unitary has eigenvalues of modulus one; $(e_{1}e_{2})^{2} = -e_{1}^{2}e_{2}^{2} = 1$ and $e_{1}e_{2} \ne \pm 1$, so $\exp(\tfrac{\theta}{2}e_{1}e_{2}) = \cosh\tfrac{\theta}{2} + \sinh\tfrac{\theta}{2}\,e_{1}e_{2}$ has eigenvalues $e^{\pm\theta/2}$ and is unitary for no positive form when $\theta \ne 0$.

\emph{(v)} This is Theorem \claimref{synthesis.twistor-form-states-only-off-the-real-slice}(iii) at $\rho = 1$, read in the convention of Theorem \claimref{synthesis.incidence-vacua-thermofield-real-capacity}, where $e_{1}$ is the time direction (the stored twistor-form theorem's $e_{0}$; its $e_{1}$ is the spatial direction called $e_{2}$ here); the new part is the identification of the midpoint with the centre of $h$. The boost of the $(e_{1}, e_{2})$ plane at imaginary rapidity $i\theta$ sends $e_{2}$ to $\cos\theta\,e_{2} + i\sin\theta\,e_{1}$, with imaginary part $\sin\theta\,e_{1}$, future timelike exactly for $0 < \theta < \pi$; by (iii) the lines are positive there (the check confirms at $(\cos\theta, \sin\theta) = (\pm\tfrac35, \tfrac45)$ and $(0, 1)$). At $\theta = \pi/2$ the point is $ie_{1}$, whose line is the $+1$ eigenplane of $J$ by (ii). For $\rho e_{2}$ the midpoint is $i\rho e_{1}$, which equals $ie_{1}$ only for $\rho = 1$.

\emph{Reading (argued, not part of the statement).} Theorem \claimref{synthesis.twistor-form-states-only-off-the-real-slice} put the states of the twistor form off the real slice, on the complex points with future timelike imaginary part, and Theorems \claimref{synthesis.incidence-vacua-thermofield-real-capacity} and \claimref{synthesis.boost-stationary-incidence-vacua-read-interval} computed thermal content at real points under the positive structure; the open problem asked which of the two the selected state reads. By (ii) and (iii) the positive structure is the twistor form together with the point $ie_{1}$ of the forward tube, so the second family of results is the first read from that point, and by (iv) a state has to choose such a point. The incidence line of the chosen point is a definite plane, available only off the real slice, which is a complex-capacity object; the marginals at real points are the thermofield doubles of the ascent mode with the boost qubit of the stored synthesis, a real-capacity object. So one state reads both capacities, the complex one through its centre and the real one through its real-point marginals, and the dependence on frame and scale noted in clause (iv) of the thermofield synthesis is the dependence on the tube point. The point $ie_{1}$ is the midpoint of the positive half strip of the imaginary boost of the unit real point, the place a Bisognano-Wichmann-type modular structure would put its vacuum; that identification is not proved here (in the Bisognano-Wichmann picture, cited from memory, vectors are analytic in the imaginary rapidity on the half strip $0 < \theta < \pi$, half of the KMS strip of width $2\pi$, and the midpoint $\theta = \pi/2$ is where the modular operator enters to the power one quarter; that the time orientation for which $L(ie_{1})$ is positive is the one in which the null translations have positive generator is argued only). Clause (iv) is also a finite-dimensional reason why a boost-KMS state cannot live on the eight-dimensional module: no positive structure there is boost-invariant, so the boost acts unitarily only on a continuum, where the result filed alongside on positive null energy (synthesis.positive-null-energy-gives-r-premises, not yet in the store) poses the corresponding premise and reaches the KMS strip of width $2\pi$, of which the positive half strip $0 < \theta < \pi$ here is the vector-valued half (the half on which the modular operator enters to the power one half), with midpoint at the power one quarter. The link between the time orientation that makes $L(ie_{1})$ positive and positive null energy is argued, not proved. Nothing here connects the tube point to the selected state of the boundary modes.
